\chapter{Nonstress Test}
\section*{What Is This Test?} A nonstress test monitors fetal heart rate and movement without inducing contractions.\section*{Alternative Names} NST; antenatal fetal monitoring; cardiotocography.\section*{Principle} A healthy fetus usually has heart-rate accelerations with movement, reflecting adequate oxygenation and nervous-system function.\section*{Why This Test Is Ordered} It monitors pregnancies with reduced movement, diabetes, hypertension, growth concerns, post-term pregnancy, or other risk.\section*{Primary vs Secondary Test} NST is a surveillance test; biophysical profile, ultrasound, or contraction stress testing may provide additional information.\section*{How This Test Is Performed} Sensors are placed on the abdomen for usually 20--40 minutes while movement is recorded.\section*{What Preparation Is Needed} Eat, drink, and take medicines as instructed; report decreased fetal movement promptly.\section*{Understanding Results} Reactive and reassuring findings differ by gestational age; nonreactive results may reflect sleep, medicines, or reduced oxygen and require clinical follow-up.\section*{Follow-up Tests} Repeat monitoring, biophysical profile, ultrasound, or delivery planning may follow.\section*{References} \begin{enumerate}\item MedlinePlus. Nonstress Test.\item American College of Obstetricians and Gynecologists. Antepartum fetal surveillance.\end{enumerate}\clearpage\section*{Understanding Results and Follow-up}
The laboratory report should identify the specimen, method, units, reference interval or decision limit, and whether the result is positive, negative, abnormal, or inconclusive. Reference intervals vary with age, sex, pregnancy, specimen type, assay, and laboratory; a result outside a range is not automatically a diagnosis, and a result within a range does not always exclude disease. Discuss unexpected results with the ordering clinician, who can decide whether repeat or confirmatory testing, treatment, or referral is appropriate. Do not change medicines or delay urgent care based on an isolated result.
\section*{Regulatory Status and Authorization Date}This chapter describes a test category, not a specific commercial product. FDA, CDSCO, EU IVDR, and MHRA approval or registration dates therefore cannot be assigned without the assay manufacturer, product name, instrument, and intended use. For laboratory-developed tests, an individual agency approval date may not exist. See the Preface for the interpretation of regulatory status.
\clearpage
